\paragraph{(I, 1.7.3) y (I, 2.2.4)}
Los enunciados dados bajo estos números pueden generalizarse, de
acuerdo con una observación de J.~Tate, como sigue.

\subsection*{1.8. Morfismos de espacios localmente anillados a esquemas afines}
\label{ega2:addendum:subsection:I.1.8-fr}

\begin{proposition}[1.8.1]
\label{ega2:addendum:I.1.8.1-fr}
Sean $(S,\sh{O}_S)$ un esquema afín y $(X,\sh{O}_X)$ un espacio
localmente anillado. Existe una biyección canónica del conjunto de
homomorfismos de anillos
\oldpage[II]{218}
$\Gamma(S,\sh{O}_S)\to\Gamma(X,\sh{O}_X)$
al conjunto de morfismos de espacios anillados
$(\psi,\theta):(X,\sh{O}_X)\to(S,\sh{O}_S)$ tales que, para todo
$x\in X$, $\theta_x^\sharp$ sea un homomorfismo local: $\sh{O}_{\psi(x)}\to\sh{O}_x$.
\end{proposition}

\begin{proof}
Observemos primero que, si $(X,\sh{O}_X)$ y $(S,\sh{O}_S)$ son dos
espacios anillados cualesquiera, un morfismo $(\psi,\theta)$ de
$(X,\sh{O}_X)$ a $(S,\sh{O}_S)$ define canónicamente un homomorfismo
de anillos
$\Gamma(\theta):\Gamma(S,\sh{O}_S)\to\Gamma(X,\sh{O}_X)$;
y, por tanto, una primera aplicación
\begin{equation}
\label{ega2:addendum:I.1.8.1.1-fr}
  \rho:\Hom((X,\sh{O}_X),(S,\sh{O}_S))
  \to\Hom(\Gamma(S,\sh{O}_S),\Gamma(X,\sh{O}_X)).
\tag{1.8.1.1}
\end{equation}
Recíprocamente, bajo las hipótesis de la proposición, pongamos
$A=\Gamma(S,\sh{O}_S)$ y consideremos un homomorfismo de anillos
$\vphi:A\to\Gamma(X,\sh{O}_X)$. Para todo $x\in X$, el conjunto de
los $f\in A$ tales que $\vphi(f)(x)=0$ es un \emph{ideal primo} de $A$, pues
$\sh{O}_x/\mathfrak m_x=\boldsymbol{k}(x)$ es un cuerpo. Es, por tanto, un
elemento de $S=\Spec(A)$, que denotaremos asimismo por
${}^a\vphi(x)$. Además, para todo $f\in A$, por definición (\textbf 0, 5.5.2),
${}^a\vphi^{-1}(D(f))=X_f$,
lo que prueba que ${}^a\vphi$ es una aplicación \emph{continua} $X\to S$.

Definamos después un homomorfismo
\[
  \widetilde{\vphi}:\sh{O}_S\to{}^a\vphi_*(\sh{O}_X)
\] de $\sh{O}_S$-módulos; para todo $f\in A$, se tiene
$\Gamma(D(f),\sh{O}_S)=A_f$ (1.3.6); para todo $s\in A$, se hace corresponder a $s/f\in A_f$ el elemento
$(\vphi(s)|X_f)(\vphi(f)|X_f)^{-1}$ de $\Gamma(X_f,\sh{O}_X)=\Gamma(D(f),{}^a\vphi_*(\sh{O}_X))$,
y se comprueba inmediatamente (pasando de $D(f)$ a $D(fg)$) que esto define un homomorfismo de $\sh{O}_S$-módulos, de donde, un morfismo de espacios
anillados $({}^a\vphi,\widetilde{\vphi})$. Con la misma notación, y
escribiendo $y={}^a\vphi(x)$, se tiene (\textbf 0, 3.7.1)
$\widetilde{\vphi}_x^\sharp(s_y/f_y)=(\vphi(s)_x)(\vphi(f)_x)^{-1}$;
Como $s_y\in\mathfrak m_y$ equivale, por definición, a
$\vphi(s)_x\in\mathfrak m_x$, se ve que $\widetilde{\vphi}_x^\sharp$ es un homomorfismo \emph{local} $\sh{O}_y\to\sh{O}_x$. Se ha
definido así una segunda aplicación
$\sigma:\Hom(\Gamma(S,\sh{O}_S),\Gamma(X,\sh{O}_X))\to\mathfrak L$,
donde $\mathfrak L$ denota el conjunto de morfismos
$(\psi,\theta):(X,\sh{O}_X)\to(S,\sh{O}_S)$ para los cuales
$\theta_x^\sharp$ es local para todo $x\in X$.

Queda probar que $\sigma$ y la restricción de $\rho$ a $\mathfrak L$
son inversas. La definición de $\widetilde{\vphi}$ da inmediatamente
$\Gamma(\widetilde{\vphi})=\vphi$, y por tanto $\rho\circ\sigma$ es la
identidad. Para ver que $\sigma\circ\rho$ es la identidad, partamos de
$(\psi,\theta)\in\mathfrak L$ y sea $\vphi=\Gamma(\theta)$. La
hipótesis sobre $\theta_x^\sharp$ proporciona, al pasar a los cuerpos
residuales, un monomorfismo
$\theta^x:\boldsymbol{k}(\psi(x))\to\boldsymbol{k}(x)$
tal que, para todo $f\in A=\Gamma(S,\sh{O}_S)$,
$\theta^x(f(\psi(x)))=\vphi(f)(x)$;
Así, $f(\psi(x))=0$ equivale a $\vphi(f)(x)=0$, lo que prueba que
${}^a\vphi=\psi$. Por otra parte, las definiciones hacen conmutativo
el diagrama
\[
  \xymatrix{
    A\ar[r]^\vphi\ar[d] &
    \Gamma(X,\sh{O}_X)\ar[d]\\
    A_{\psi(x)}\ar[r]_{\theta_x^\sharp} &
    \sh{O}_x
  }
\]
y también es conmutativo el diagrama análogo en el que
$\theta_x^\sharp$ se reemplaza por $\widetilde{\vphi}_x^\sharp$. De
(\textbf 0, 1.2.4) se sigue que
$\widetilde{\vphi}_x^\sharp=\theta_x^\sharp$, y, por consiguiente,
$\widetilde{\vphi}=\theta$.
\end{proof}

\begin{env}[1.8.2]
\label{ega2:addendum:I.1.8.2-fr}
Cuando $(X,\sh{O}_X)$ e $(Y,\sh{O}_Y)$ son espacios \emph{localmente}
anillados, consideraremos los morfismos
$(\psi,\theta):(X,\sh{O}_X)\to(Y,\sh{O}_Y)$ tales que, para todo
$x\in X$, $\theta_x^\sharp$ sea un homomorfismo \emph{local}: $\sh{O}_{\psi(x)}\to\sh{O}_x$.

Siempre que en lo sucesivo hablemos de un
\oldpage[II]{219}
\emph{morfismo de espacios localmente anillados}, entenderemos siempre
un morfismo de este tipo. Con esta definición de los morfismos, los
espacios localmente anillados forman una \emph{categoría}. Para dos objetos
$X,Y$ de esta categoría, $\Hom(X,Y)$ denotará, por tanto, el conjunto
de morfismos de espacios localmente anillados de $X$ a $Y$ (el
conjunto denotado por $\mathfrak L$ en \sref{ega2:addendum:I.1.8.1-fr}). Cuando queramos
considerar el conjunto de \emph{morfismos de espacios anillados} de $X$ a
$Y$, escribiremos $\Hom_{\mathit{an}}(X,Y)$ para evitar toda
ambigüedad. La aplicación \sref{ega2:addendum:I.1.8.1.1-fr} se escribe entonces
\begin{equation}
\label{ega2:addendum:I.1.8.2.1-fr}
  \rho:\Hom_{\mathit{an}}(X,Y)
  \to\Hom(\Gamma(Y,\sh{O}_Y),\Gamma(X,\sh{O}_X))
\tag{1.8.2.1}
\end{equation}
y su restricción
\begin{equation}
\label{ega2:addendum:I.1.8.2.2-fr}
  \rho':\Hom(X,Y)
  \to\Hom(\Gamma(Y,\sh{O}_Y),\Gamma(X,\sh{O}_X))
\tag{1.8.2.2}
\end{equation}
es \emph{funtorial} en $X$ e $Y$ en la categoría de los espacios localmente
anillados.
\end{env}

\begin{corollary}[1.8.3]
\label{ega2:addendum:I.1.8.3-fr}
Sea $(Y,\sh{O}_Y)$ un espacio localmente anillado. Para que $Y$ sea un
esquema afín, es necesario y suficiente que, para todo espacio
localmente anillado $(X,\sh{O}_X)$, la aplicación \sref{ega2:addendum:I.1.8.2.2-fr} sea
biyectiva.
\end{corollary}

\begin{proof}
La Proposición~\sref{ega2:addendum:I.1.8.1-fr} muestra que la condición es necesaria.
Recíprocamente, se supone que se cumple y se pone
$A=\Gamma(Y,\sh{O}_Y)$. La hipótesis y \sref{ega2:addendum:I.1.8.1-fr} muestran que los
funtores
$X\to\Hom(X,Y)$ y $X\to\Hom(X,\Spec(A))$
de los espacios localmente anillados a los conjuntos son \emph{isomorfos}.
Esto implica la existencia de un isomorfismo canónico
$Y\to\Spec(A)$ (cf.~\textbf 0, 8).
\end{proof}

\begin{env}[1.8.4]
\label{ega2:addendum:I.1.8.4-fr}
Sea $S=\Spec(A)$ un esquema afín. Denotemos por $(S',A')$ el espacio
anillado cuyo espacio subyacente \emph{consta de un punto} y cuyo haz
estructural $A'$ es el haz (necesariamente simple) sobre $S'$ definido
por el anillo $A$. Sea $\pi:S\to S'$ la única aplicación de $S$ en $S'$; observemos, por otra parte, que para todo
abierto $U$ de $S$ se tiene una aplicación canónica
$\Gamma(S',A')=\Gamma(S,\sh{O}_S)\to\Gamma(U,\sh{O}_S)$
que define un \emph{$\pi$-morfismo} $\iota:A'\to\sh{O}_S$ de haces de anillos.
Se obtiene así un \emph{morfismo} canónico \emph{de espacios anillados}
$i=(\pi,\iota):(S,\sh{O}_S)\to(S',A')$.
Para todo $A$-módulo $M$, denotemos por $M'$ el haz simple sobre $S'$
definido por $M$, que es evidentemente un $A'$-módulo. Está claro que
$i_*(\widetilde M)=M'$ (1.3.7).
\end{env}

\begin{lemma}[1.8.5]
\label{ega2:addendum:I.1.8.5-fr}
Con la notación de \sref{ega2:addendum:I.1.8.4-fr}, para todo $A$-módulo $M$, el
$\sh{O}_S$-homomorfismo canónico y funtorial \textup{(\textbf{0}, 4.4.3.3)}
\begin{equation}
\label{ega2:addendum:I.1.8.5.1-fr}
  i^*(i_*(\widetilde M))\to\widetilde M
\tag{1.8.5.1}
\end{equation}
es un isomorfismo.
\end{lemma}

\begin{proof}
Ambos miembros de \sref{ega2:addendum:I.1.8.5.1-fr} son exactos por la derecha
(siendo evidentemente exacto el funtor
$M\to i_*(\widetilde M)$) y conmutan con las sumas directas.
Considerando $M$ como el conúcleo de un homomorfismo
$A^{(I)}\to A^{(J)}$, se reduce uno al caso $M=A$, en el cual la
afirmación es inmediata.
\end{proof}

\begin{corollary}[1.8.6]
\label{ega2:addendum:I.1.8.6-fr}
Sea $(X,\sh{O}_X)$ un espacio anillado y sea $u:X\to S$ un morfismo de
\oldpage[II]{220}
espacios anillados. Para todo $A$-módulo $M$, existe, con la notación
de \sref{ega2:addendum:I.1.8.4-fr}, un isomorfismo canónico y funtorial de
$\sh{O}_X$-módulos
\begin{equation}
\label{ega2:addendum:I.1.8.6.1-fr}
  u^*(\widetilde M)\xrightarrow{\sim} u^*(i^*(M'))
\tag{1.8.6.1}
\end{equation}
\end{corollary}

